\section{Local score calculus and the convex-ridge obstruction}
\label{app:local-ridge}

\subsection{Second-score identities and the non-Gaussian modulus}
\label{app:local-moduli}

This subsection proves the score identities and the explicit bounds
\eqref{eq:local-Kng-explicit}--\eqref{eq:local-Kng-covariance}.  We work in
the optimizer-whitened coordinates of \Cref{sec:local-gaussian-core}.  Assume
throughout that \(\widehat V\in C^2(\R^d)\) and
\[
  \alphastar\Id\preceq\nabla^2\widehat V(x)
  \preceq\betastar\Id,\qquad x\in\R^d.
\]
Boundedness of the Hessian implies at most quadratic growth of
\(\widehat V\) and at most linear growth of its gradient.  It therefore
justifies all Gaussian density differentiations below, locally uniformly on
sets whose covariance eigenvalues are bounded away from zero and infinity.
No classical third or fourth derivative is used.  The formulas also apply
entrywise to vector- and matrix-valued integrands.

Fix \(a=(m,C)\), a tangent vector \(\zeta=(u,U)\), and set
\begin{equation}
  v:=C^{-1/2}u,\qquad
  S:=C^{-1/2}UC^{-1/2},\qquad
  t^2:=\norm v^2+\frac12\norm S_{\Frob}^2.
  \label{eq:app-local-whitened-tangent}
\end{equation}
If \(Z\sim\N(0,\Id)\) and \(X=m+C^{1/2}Z\), define
\begin{align}
  q_\zeta(Z)
  &:=
  v^\top Z+\frac12\bigl(Z^\top SZ-\Tr S\bigr),
  \label{eq:app-local-first-score}\\
  p_\zeta(Z)
  &:=
  \frac12\Tr(S^2)-\norm v^2-2v^\top SZ-Z^\top S^2Z,
  \qquad
  \chi_\zeta:=q_\zeta^2+p_\zeta.
  \label{eq:app-local-second-score}
\end{align}

\begin{lemma}[First and second Gaussian score identities]
\label{lem:app-local-second-score}
Let \(\phi:\R^d\to\R\) be measurable and of polynomial growth, and put
\(\Phi(m,C):=\E_{\N(m,C)}[\phi]\).  Then
\begin{equation}
  D\Phi(a)[\zeta]=\E[\phi(X)q_\zeta(Z)],
  \qquad
  D^2\Phi(a)[\zeta,\zeta]
  =\E[\phi(X)\chi_\zeta(Z)].
  \label{eq:app-local-score-identities}
\end{equation}
Moreover, \(q_\zeta\) belongs to the first two Hermite chaoses, whereas
\(\chi_\zeta\) belongs to chaoses two, three, and four.  In particular,
\begin{equation}
  \E q_\zeta=\E\chi_\zeta=0,\qquad
  \E[Z\chi_\zeta]=0,\qquad
  \E q_\zeta^2=t^2,\qquad
  \norm{\chi_\zeta}_{L^2}\leq\sqrt6\,t^2.
  \label{eq:app-local-score-moments}
\end{equation}
The constant \(\sqrt6\) is sharp.
\end{lemma}

\begin{proof}
Let \(m_s=m+su\) and \(C_s=C+sU\), which remains positive definite for
small \(\abs s\).  Differentiating the log-density of
\(\N(m_s,C_s)\) at the fixed point \(x=m+C^{1/2}Z\) gives
\[
  \left.\partial_s\log\rho_{m_s,C_s}(x)\right|_{s=0}=q_\zeta(Z),
  \qquad
  \left.\partial_s^2\log\rho_{m_s,C_s}(x)\right|_{s=0}=p_\zeta(Z).
\]
The identity
\(\partial_s^2\rho_s=\rho_s[(\partial_s\log\rho_s)^2+
\partial_s^2\log\rho_s]\) proves
\eqref{eq:app-local-score-identities}.

For the chaos calculation, write
\[
  \ell:=v^\top Z,\qquad
  Q:=Z^\top SZ-\Tr S,\qquad
  R_2:=Z^\top S^2Z-\Tr S^2.
\]
Then \(q_\zeta=\ell+Q/2\), and orthogonality of the first and second
chaoses gives
\(\E q_\zeta^2=\norm v^2+\norm S_{\Frob}^2/2=t^2\).
The Gaussian product formula yields
\[
  (Q^2)_0=2\norm S_{\Frob}^2,\qquad
  (Q^2)_2=4R_2,\qquad
  \ell Q=2v^\top SZ+G_3,
\]
where \(G_3\) is pure chaos three.  Substitution in
\(\chi_\zeta=q_\zeta^2+p_\zeta\) cancels its chaos-zero and chaos-one
parts and leaves the mutually orthogonal components
\[
  (\chi_\zeta)_2=\ell^2-\norm v^2,\qquad
  (\chi_\zeta)_3=G_3,\qquad
  (\chi_\zeta)_4=\frac14(Q^2)_4.
\]
This proves the first three assertions in
\eqref{eq:app-local-score-moments}.

For the \(L^2\) norm, abbreviate
\[
  a_0:=\norm v^2,\quad
  b_0:=\norm S_{\Frob}^2,\quad
  c_0:=v^\top S^2v,\quad
  d_4:=\Tr(S^4).
\]
Wick's formula and the preceding orthogonal decomposition give
\begin{equation}
  \norm{\chi_\zeta}_{L^2}^2
  =
  2a_0^2+2a_0b_0+4c_0+\frac12b_0^2+d_4.
  \label{eq:app-local-score-exact-L2}
\end{equation}
Since \(c_0\leq a_0b_0\) and \(d_4\leq b_0^2\),
\[
  \norm{\chi_\zeta}_{L^2}^2
  \leq
  2a_0^2+6a_0b_0+\frac32b_0^2
  \leq6\left(a_0+\frac{b_0}{2}\right)^2
  =6t^4.
\]
Equality holds when \(v=0\) and \(S\) has rank one, proving sharpness.
\end{proof}

The Ornstein--Uhlenbeck generator
\(\mathcal L_{\mathrm{OU}}=\Delta_z-z^\top\nabla_z\) acts as multiplication
by \(-k\) on the \(k\)-th Hermite chaos.  Thus, for a centered
\(\psi=\sum_{k\geq1}\psi_k\),
\begin{equation}
  \E\norm{\nabla(-\mathcal L_{\mathrm{OU}})^{-1}\psi}^2
  =
  \sum_{k\geq1}\frac1k\norm{\psi_k}_{L^2}^2.
  \label{eq:app-local-OU-inverse}
\end{equation}
In particular, the right-hand side is at most
\(\norm\psi_{L^2}^2/2\) when \(\psi\) has no chaos-one component.  This
factor \(1/2\) is applicable to \(\chi_\zeta\), but not to a general
centered score.

\begin{lemma}[Explicit non-Gaussian second-differential modulus]
\label{lem:app-local-modulus-bound}
Let \(H=F-\FGauss\) be the non-Gaussian part of the whitened
Fisher--Rao field.  On the star hull \(\mathcal H_\delta\),
\[
  H_m(a)=C(g(a)+m),\qquad
  H_C(a)=-C(A(a)-\Id)C.
\]
Then
\begin{equation}
  \Kng(\delta)
  :=
  \sup_{\substack{a\in\mathcal H_\delta\\\zeta\neq0}}
  \frac{\norm{D^2H(a)[\zeta,\zeta]}_\star}
       {\norm{\zeta}_\star^2}
  \leq
  \left(
    K_{\mathrm{ng},m}(\delta)^2
    +\frac12K_{\mathrm{ng},C}(\delta)^2
  \right)^{1/2},
  \label{eq:app-local-modulus-pair}
\end{equation}
Here the two block constants are
\begin{align}
  K_{\mathrm{ng},m}(\delta)
  &:=
  \frac{U_\delta^{3/2}}{\ell_\delta^2}
  \left(
    \sqrt3\,\overline\Sigma_\delta+\Sigma_\delta
    +\sqrt{\Sigma_\delta^2+2\mathcal B_\delta^2}
  \right),
  \label{eq:app-local-Kng-mean}\\
  K_{\mathrm{ng},C}(\delta)
  &:=
  \frac{U_\delta^2}{\ell_\delta^2}
  \left(
    (4\sqrt2+\sqrt6)\Sigma_\delta+4\mathcal B_\delta
  \right).
  \label{eq:app-local-Kng-covariance}
\end{align}
These are the quantities stated in
\eqref{eq:local-Kng-mean}--\eqref{eq:local-Kng-covariance}.
\end{lemma}

\begin{proof}
Fix \(a\in\mathcal H_\delta\).  The covariance band gives
\(\ell_\delta\Id\preceq C\preceq U_\delta\Id\).  Put
\[
  x:=\norm v,\qquad
  y:=\frac{\norm S_{\Frob}}{\sqrt2},
  \qquad t^2=x^2+y^2.
\]
Then
\begin{equation}
  \norm U_{\op}\leq\sqrt2U_\delta y,\qquad
  \norm U_{\Frob}\leq\sqrt2U_\delta y,\qquad
  t\leq\ell_\delta^{-1}\norm\zeta_\star.
  \label{eq:app-local-band-conversion}
\end{equation}
The final inequality uses \(\ell_\delta<1\).

For the mean block, Gaussian integration by parts through
\(\psi_1=(-\mathcal L_{\mathrm{OU}})^{-1}q_\zeta\), for which
\(\nabla\psi_1=v+SZ/2\), gives
\begin{equation}
  D(g+m)[\zeta]
  =
  -\E\!\left[
    \bigl(\nabla^2\widehat V(X)-\Id\bigr)
    C^{1/2}\left(v+\frac12SZ\right)
  \right].
  \label{eq:app-local-first-g-derivative}
\end{equation}
Split
\(\nabla^2\widehat V-\Id=(\nabla^2\widehat V-A)+(A-\Id)\).
The centered Hessian annihilates the constant \(v\), while the deterministic
part annihilates \(SZ\).  By the definitions of
\(\mathcal B_\delta\) and \(\Sigma_\delta\),
\begin{equation}
  \norm{D(g+m)[\zeta]}
  \leq
  \sqrt{U_\delta}
  \left(\mathcal B_\delta x+\frac{\Sigma_\delta}{\sqrt2}y\right).
  \label{eq:app-local-first-g-bound}
\end{equation}

For the second derivative, let
\(\psi_2=(-\mathcal L_{\mathrm{OU}})^{-1}\chi_\zeta\).  By
\eqref{eq:app-local-OU-inverse} and
\eqref{eq:app-local-score-moments},
\[
  \E\norm{\nabla\psi_2}^2\leq3t^4,\qquad
  \E\nabla\psi_2=0.
\]
Another Gaussian integration by parts therefore gives
\[
  D^2g[\zeta,\zeta]
  =
  -\E\!\left[
    \bigl(\nabla^2\widehat V(X)-A(a)\bigr)
    C^{1/2}\nabla\psi_2
  \right].
\]
This expectation can be bounded either with the Hessian dispersion
\(\Sigma_\delta\) or with the pointwise operator bound
\(\norm{\nabla^2\widehat V-A}_{\op}\leq\betastar-\alphastar\).
Consequently,
\begin{equation}
  \norm{D^2g[\zeta,\zeta]}
  \leq
  \sqrt{3U_\delta}\,\overline\Sigma_\delta\,t^2.
  \label{eq:app-local-second-g-bound}
\end{equation}
This is an operator-norm alternative, not the false pointwise assertion
\(\norm{\nabla^2\widehat V-A}_{\Frob}
\leq\betastar-\alphastar\).

Since \(C\) is affine in the covariance coordinate,
\[
  D^2H_m[\zeta,\zeta]
  =
  2U\,D(g+m)[\zeta]+C\,D^2g[\zeta,\zeta].
\]
Combining \eqref{eq:app-local-first-g-bound} and
\eqref{eq:app-local-second-g-bound} with
\eqref{eq:app-local-band-conversion} yields
\[
  \norm{D^2H_m[\zeta,\zeta]}
  \leq
  U_\delta^{3/2}
  \left(
    2\sqrt2\,\mathcal B_\delta xy
    +2\Sigma_\delta y^2
    +\sqrt3\,\overline\Sigma_\delta t^2
  \right).
\]
The largest eigenvalue of the first two terms as a quadratic form in
\((x,y)\) is
\(\Sigma_\delta+\sqrt{\Sigma_\delta^2+2\mathcal B_\delta^2}\).
Using \(t\leq\ell_\delta^{-1}\norm\zeta_\star\) proves
\begin{equation}
  \norm{D^2H_m[\zeta,\zeta]}
  \leq
  K_{\mathrm{ng},m}(\delta)\norm\zeta_\star^2.
  \label{eq:app-local-mean-modulus}
\end{equation}

For the covariance block, write \(W=A-\Id\).  The score identities imply
\begin{equation}
  \norm{DW[\zeta]}_{\Frob}\leq\Sigma_\delta t,\qquad
  \norm{D^2W[\zeta,\zeta]}_{\Frob}
  \leq\sqrt6\,\Sigma_\delta t^2.
  \label{eq:app-local-A-score-bounds}
\end{equation}
Indeed, \(DA[\zeta]=\E[(\nabla^2\widehat V-A)q_\zeta]\) and
\(D^2A[\zeta,\zeta]
=\E[(\nabla^2\widehat V-A)\chi_\zeta]\).
The product rule gives
\[
  D^2(CWC)[\zeta,\zeta]
  =
  2U(DW)C+2C(DW)U+2UWU+C(D^2W)C.
\]
Using \(\norm W_{\op}\leq\mathcal B_\delta\) and
\eqref{eq:app-local-band-conversion}--\eqref{eq:app-local-A-score-bounds},
\begin{align*}
  \norm{D^2H_C[\zeta,\zeta]}_{\Frob}
  &\leq
  U_\delta^2\left(
    4\sqrt2\,\Sigma_\delta ty
    +4\mathcal B_\delta y^2
    +\sqrt6\,\Sigma_\delta t^2
  \right)\\
  &\leq
  K_{\mathrm{ng},C}(\delta)\norm\zeta_\star^2.
\end{align*}
Finally,
\(\norm{D^2H}_\star^2
=\norm{D^2H_m}^2+\norm{D^2H_C}_{\Frob}^2/2\), which proves
\eqref{eq:app-local-modulus-pair}.  Each constant is proportional to
\(\mathcal B_\delta\), \(\Sigma_\delta\), or
\(\overline\Sigma_\delta\), so the modulus vanishes for a Gaussian target.
\end{proof}

\subsection{The convex-ridge construction}
\label{app:local-ridge-construction}

All constants implicit in \(O(\cdot)\) below are numerical and independent
of \(K\) and the dimension.  Let \(p\in\mathbb N\) tend to infinity and set
\begin{equation}
  K:=(8p)^{1/4},\qquad
  d_K:=p+1,\qquad
  p=\frac{K^4}{8}.
  \label{eq:app-ridge-dimension-scale}
\end{equation}
Write \(z=(t,y)\in\R\times\R^p\), \(r=\norm y\), and use the regularized
radius
\begin{equation}
  \rho(y):=\sqrt{p+r^2},\qquad
  \rho_0:=\sqrt{2p}=\frac{K^2}{2},\qquad
  c:=\rho_0+\frac12.
  \label{eq:app-ridge-radius}
\end{equation}
The additive \(p\) in \(\rho\) is essential: it keeps \(D^2\rho\) and
\(D^3\rho\) uniformly controlled at \(y=0\).

Define a convex \(C^{2,1}\) profile \(F_K:\R\to\R\), up to an irrelevant
additive constant, by
\begin{equation}
  F_K''(x):=(K-\abs{x-c})_+,\qquad
  F_K'(-\infty)=0.
  \label{eq:app-ridge-profile}
\end{equation}
Its total curvature mass is \(K^2\), and
\begin{equation}
  0\leq F_K''\leq K,\qquad
  0\leq F_K'\leq K^2,\qquad
  \Lip(F_K'')\leq1.
  \label{eq:app-ridge-profile-bounds}
\end{equation}
For \(s\geq0\), set
\begin{equation}
  U_s(t,y):=F_K(\rho(y)+st),\qquad
  x:=\rho(y)+st,\qquad
  q:=\nabla x=\left(s,\frac y\rho\right).
  \label{eq:app-ridge-ridge}
\end{equation}
Since
\[
  D^2\rho(y)=\frac{\Id_p}{\rho}-\frac{yy^\top}{\rho^3}\succeq0,
\]
the ridge Hessian is
\begin{equation}
  \mathcal A_s:=D^2U_s
  =
  F_K''(x)qq^\top+
  F_K'(x)
  \begin{pmatrix}
    0&0\\
    0&D^2\rho(y)
  \end{pmatrix}
  \succeq0.
  \label{eq:app-ridge-ridge-Hessian}
\end{equation}

Let \(Z=(\Theta,Y)\sim\N(0,\Id_{d_K})\), put
\(R=\norm Y\), and continue to write
\(\rho=\sqrt{p+R^2}\).  Define \(\xi:=\rho-\rho_0\).

\begin{lemma}[Dimension-free radial fluctuations]
\label{lem:app-ridge-radial-concentration}
For every fixed \(q<\infty\), there is \(C_q<\infty\) such that
\begin{equation}
  \norm{\xi}_{L^q}\leq C_q,
  \qquad
  \E\xi=-\frac{\E\xi^2}{2\rho_0}=O(K^{-2}).
  \label{eq:app-ridge-radial-moments}
\end{equation}
Uniformly for \(0\leq s\leq2K^{-1/2}\), the random variable
\[
  D_s:=\xi+s\Theta-\frac12
\]
has dimension-free sub-Gaussian tails.  In particular, for every fixed
\(q<\infty\), there are \(c_q,C_q>0\) such that
\begin{equation}
  \sup_{0\leq s\leq2K^{-1/2}}
  \E\!\left[(1+\abs{D_s})^q\,
    \mathbf 1_{\{\abs{D_s}>K\}}\right]
  \leq C_qe^{-c_qK^2}.
  \label{eq:app-ridge-radial-tail}
\end{equation}
\end{lemma}

\begin{proof}
The identity
\[
  \xi=\frac{R^2-p}{\rho+\rho_0}
\]
and \(\rho+\rho_0\geq\rho_0\asymp\sqrt p\) give
\(\norm\xi_{L^q}\leq
\norm{R^2-p}_{L^q}/\rho_0\leq C_q\); the last inequality is the standard
fixed-moment bound for a centered \(\chi_p^2\) variable.  Moreover,
\(\rho^2-\rho_0^2=2\rho_0\xi+\xi^2\) has mean zero, which proves the exact
mean identity in \eqref{eq:app-ridge-radial-moments}.

The map \(y\mapsto\sqrt{p+\norm y^2}\) is \(1\)-Lipschitz.  Gaussian
concentration therefore makes \(\rho-\E\rho\) sub-Gaussian with a numerical
variance proxy.  Since \(\E\xi=O(K^{-2})\) and
\(s\Theta\) has variance at most \(4/K\), the same is true uniformly for
\(D_s\).  Integrating the sub-Gaussian tail, or applying Cauchy--Schwarz
against its fixed polynomial moments, proves
\eqref{eq:app-ridge-radial-tail}.
\end{proof}

The tilt is chosen to make the averaged ridge Hessian isotropic.  Define
\begin{align}
  \overline K(s)
  &:=
  \E F_K''(\rho+s\Theta),
  \label{eq:app-ridge-Kbar}\\
  G(s)
  &:=
  \frac1p\E\!\left[
    F_K''(\rho+s\Theta)\frac{R^2}{\rho^2}
    +
    F_K'(\rho+s\Theta)
    \left(\frac p\rho-\frac{R^2}{\rho^3}\right)
  \right].
  \label{eq:app-ridge-G}
\end{align}

\begin{lemma}[Averaged profile and exact isotropic tuning]
\label{lem:app-ridge-tuning}
For every \(s\geq0\),
\begin{equation}
  \E\mathcal A_s
  =
  \diag\bigl(s^2\overline K(s),G(s)\Id_p\bigr).
  \label{eq:app-ridge-average-Hessian}
\end{equation}
Uniformly for \(0\leq s\leq2K^{-1/2}\),
\begin{equation}
  \overline K(s)=K+O(1),\qquad
  \E F_K'(\rho+s\Theta)=\rho_0-\frac K2+O(1),\qquad
  G(s)=1-\frac1K+O(K^{-2}).
  \label{eq:app-ridge-profile-estimates}
\end{equation}
For all sufficiently large \(K\), there is
\(s_K\in(0,2K^{-1/2})\) such that
\begin{equation}
  s_K^2\overline K(s_K)=G(s_K)=:a_K.
  \label{eq:app-ridge-tuning-equation}
\end{equation}
For this choice,
\begin{equation}
  s_K^2=K^{-1}+O(K^{-2}),\qquad
  a_K=1-K^{-1}+O(K^{-2}),\qquad
  \alpha_K:=1-a_K=K^{-1}+O(K^{-2}),
  \label{eq:app-ridge-tuned-scaling}
\end{equation}
and \(\E\mathcal A_{s_K}=a_K\Id_{d_K}\).
\end{lemma}

\begin{proof}
The \((t,t)\) entry of \(\E\mathcal A_s\) is
\(s^2\overline K(s)\).  Every mixed entry is odd under one of the
reflections \(Y_i\mapsto-Y_i\).  Rotational symmetry makes the \(yy\)
block a scalar multiple of \(\Id_p\), and averaging its diagonal entries
using \eqref{eq:app-ridge-ridge-Hessian} gives exactly \(G(s)\).  This
proves \eqref{eq:app-ridge-average-Hessian}.

Since \(\rho+s\Theta=c+D_s\), on \(\{\abs{D_s}\leq K\}\) the profile has
the exact form
\begin{equation}
  F_K''(c+D_s)=K-\abs{D_s},\qquad
  F_K'(c+D_s)
  =
  \frac{K^2}{2}+KD_s-\frac12D_s\abs{D_s}.
  \label{eq:app-ridge-local-profile}
\end{equation}
The tail contribution outside this event is exponentially small by
\eqref{eq:app-ridge-radial-tail}.  Since
\(\E\abs{D_s}=O(1)\), \(\E[D_s\abs{D_s}]=O(1)\), and
\(\E D_s=-1/2+O(K^{-2})\), formula
\eqref{eq:app-ridge-local-profile} gives the first two estimates in
\eqref{eq:app-ridge-profile-estimates}.

For the last estimate, decompose \(G=(\mathrm I)+(\mathrm{II})
-(\mathrm{III})\), where
\[
  (\mathrm I)=\frac1p\E\!\left[F_K''\frac{R^2}{\rho^2}\right],
  \quad
  (\mathrm{II})=\E\!\left[\frac{F_K'}{\rho}\right],
  \quad
  (\mathrm{III})=\frac1p\E\!\left[F_K'\frac{R^2}{\rho^3}\right].
\]
The profile bounds imply
\((\mathrm I)=O(K/p)=O(K^{-3})\).  Also,
\[
  \left\|\frac1\rho-\frac1{\rho_0}\right\|_{L^2}
  =
  \left\|\frac{\xi}{\rho\rho_0}\right\|_{L^2}
  \leq\frac{\norm\xi_{L^2}}{\sqrt2\,p}
  =O(K^{-4}),
\]
and \(F_K'\leq K^2\), so
\[
  (\mathrm{II})
  =
  \frac{\E F_K'}{\rho_0}+O(K^{-2}).
\]
Finally,
\(\sup_{r\geq0}r^2(p+r^2)^{-3/2}\leq p^{-1/2}\), whence
\((\mathrm{III})\leq K^2p^{-3/2}=O(K^{-4})\).
Using \(\rho_0=K^2/2\) proves
\(G(s)=1-K^{-1}+O(K^{-2})\).

The function
\(\Phi_K(s):=s^2\overline K(s)-G(s)\) is continuous by dominated
convergence.  The preceding estimates give
\[
  \Phi_K(0)=-1+O(K^{-1})<0,\qquad
  \Phi_K(2K^{-1/2})=3+O(K^{-1})>0.
\]
The intermediate value theorem produces \(s_K\).
At a zero, \(a_K=G(s_K)=1-K^{-1}+O(K^{-2})\), and division by
\(\overline K(s_K)=K+O(1)\) gives the remaining estimates in
\eqref{eq:app-ridge-tuned-scaling}.
\end{proof}

Fix this tuned \(s=s_K\), let
\begin{equation}
  b_K:=\E[\nabla U_s(Z)],
  \qquad
  V_K(z):=\frac{\alpha_K}{2}\norm z^2+U_s(z)-b_K^\top z,
  \qquad
  B_K(z):=\nabla^2V_K(z)=\alpha_K\Id+\mathcal A_s(z).
  \label{eq:app-ridge-potential}
\end{equation}
The expectation defining \(b_K\) is finite because
\(\norm{\nabla U_s}\leq\sqrt2K^2\).

\begin{lemma}[Exact optimizer whitening and ellipticity]
\label{lem:app-ridge-whitening-ellipticity}
The standard Gaussian is the unique optimizing Gaussian for \(V_K\), and
\begin{equation}
  \E[\nabla V_K(Z)]=0,\qquad
  \E[B_K(Z)]=\Id_{d_K}.
  \label{eq:app-ridge-exact-whitening}
\end{equation}
If
\[
  L_K:=\sup_{z\in\R^{d_K}}\norm{\mathcal A_s(z)}_{\op},
  \qquad
  \beta_K:=\alpha_K+L_K,
\]
then, for all sufficiently large \(K\),
\begin{equation}
  \frac K2\leq L_K\leq3K,\qquad
  \alpha_K\Id\preceq B_K(z)\preceq\beta_K\Id,
  \qquad
  \frac{\beta_K}{\alpha_K}=\Theta(K^2).
  \label{eq:app-ridge-ellipticity}
\end{equation}
\end{lemma}

\begin{proof}
The definition of \(b_K\) and \(\E Z=0\) give the first equality in
\eqref{eq:app-ridge-exact-whitening}.  The second follows from
\(\E\mathcal A_s=a_K\Id\) and \(\alpha_K+a_K=1\).

These first-order conditions identify the optimizer, not merely a stationary
point.  Parameterize a Gaussian by \(\N(\mu,P^2)\) with \(P\in\SPD(d_K)\).
Up to a constant, its reverse-KL objective is
\[
  (\mu,P)\longmapsto
  \E[V_K(\mu+PZ)]-\log\det P.
\]
The first term is \(\alpha_K\)-strongly convex in \((\mu,P)\), because
\(\E\norm{\delta\mu+\delta PZ}^2
=\norm{\delta\mu}^2+\norm{\delta P}_{\Frob}^2\), and
\(-\log\det P\) is convex.  At \((0,\Id)\), the derivative in \(\mu\)
vanishes by \eqref{eq:app-ridge-exact-whitening}; the derivative in \(P\)
is
\(\E[\nabla V_K(Z)Z^\top]-\Id
=\E[B_K(Z)]-\Id=0\), by Gaussian integration by parts.  Strong convexity
therefore makes \((0,\Id)\) the unique minimizer.

For the upper Hessian bound, \(\rho\geq\sqrt p=K^2/\sqrt8\) and
\(s^2\leq4/K\) imply, for large \(K\),
\[
  \norm q^2=s^2+\frac{r^2}{\rho^2}\leq2,\qquad
  \norm{D^2\rho}_{\op}\leq\frac1\rho\leq\frac{\sqrt8}{K^2}.
\]
Equations \eqref{eq:app-ridge-profile-bounds} and
\eqref{eq:app-ridge-ridge-Hessian} yield
\[
  \norm{\mathcal A_s}_{\op}
  \leq
  F_K''\norm q^2+F_K'\norm{D^2\rho}_{\op}
  \leq2K+\sqrt8\leq3K.
\]
For the lower bound, take \(r^2=p\), so \(\rho=\rho_0\), and
\(t=1/(2s)\), so \(x=c\).  Then \(F_K''(x)=K\) and
\(\norm q^2=s^2+1/2\geq1/2\); the rank-one first term in
\eqref{eq:app-ridge-ridge-Hessian} has an eigenvalue at least \(K/2\).
This proves the bounds on \(L_K\).  Positive semidefiniteness of
\(\mathcal A_s\) gives the lower curvature bound, and it is attained by
taking \(t\) sufficiently negative that \(F_K'=F_K''=0\).
Finally, \(\alpha_K=\Theta(K^{-1})\) and \(L_K=\Theta(K)\) prove the
condition-number estimate.
\end{proof}

\begin{lemma}[A dimension-free Lipschitz Hessian]
\label{lem:app-ridge-Hessian-Lipschitz}
The tent-profile potential \(V_K\) belongs to \(C^{2,1}(\R^{d_K})\) and
satisfies
\begin{equation}
  \norm{\nabla^2V_K(z)-\nabla^2V_K(z')}_{\op}
  \leq64\norm{z-z'},
  \qquad z,z'\in\R^{d_K}.
  \label{eq:app-ridge-Hessian-Lipschitz}
\end{equation}
\end{lemma}

\begin{proof}
At points where the Lipschitz function \(F_K''\) is differentiable,
\(\abs{F_K'''}\leq1\).  For unit vectors \(u,v,w\), differentiating
\eqref{eq:app-ridge-ridge-Hessian} gives
\begin{align}
  D^3U_s[u,v,w]
  ={}&
  F_K'''(x)q[u]q[v]q[w]\nonumber\\
  &+
  F_K''(x)\bigl(
    q[u]D^2\rho[v,w]+q[v]D^2\rho[u,w]
    +q[w]D^2\rho[u,v]\bigr)\nonumber\\
  &+
  F_K'(x)D^3\rho[u,v,w].
  \label{eq:app-ridge-third-ridge}
\end{align}
Here derivatives of \(\rho\) ignore the distinguished coordinate.  Direct
differentiation in the \(y\) variables gives
\begin{align*}
  D^3\rho[u,v,w]
  ={}&
  -\frac{
    (u^\top v)(y^\top w)+(u^\top w)(y^\top v)
    +(v^\top w)(y^\top u)}
  {\rho^3}\\
  &+
  \frac{3(y^\top u)(y^\top v)(y^\top w)}{\rho^5}.
\end{align*}
Consequently,
\[
  \sup_{\norm u=\norm v=\norm w=1}
  \abs{D^3\rho[u,v,w]}
  \leq\frac{3r}{\rho^3}+\frac{3r^3}{\rho^5}
  \leq\frac6{\rho^2}
  \leq\frac6p=\frac{48}{K^4}.
\]
Together with
\(\norm q\leq\sqrt2\),
\(\norm{D^2\rho}_{\op}\leq\sqrt8/K^2\), and
\eqref{eq:app-ridge-profile-bounds}, this gives almost everywhere
\[
  \sup_{\norm u=\norm v=\norm w=1}
  \abs{D^3U_s[u,v,w]}
  \leq
  2\sqrt2+\frac{12}{K}+\frac{48}{K^2}
  \leq64.
\]

To pass from the almost-everywhere derivative bound to
\eqref{eq:app-ridge-Hessian-Lipschitz}, fix unit \(u,v\) and consider
\(f_{u,v}(z)=u^\top\mathcal A_s(z)v\).  This function is locally Lipschitz,
and its weak gradient has norm at most \(64\) almost everywhere.
Convolving \(f_{u,v}\) with a standard mollifier gives smooth
\(64\)-Lipschitz functions; letting the mollifier radius tend to zero
preserves the bound.  Taking the supremum over unit \(u,v\) proves the
operator-norm estimate.  The quadratic and linear terms in
\eqref{eq:app-ridge-potential} do not change the Hessian Lipschitz
seminorm.
\end{proof}

Define the first- and second-Hermite operators of the constructed potential
by
\begin{align}
  \mathcal T_K[X]
  &:=
  \E\!\left[\Tr(XB_K(Z))Z\right],
  \label{eq:app-ridge-first-Hermite-operator}\\
  \mathcal H_K[X]
  &:=
  \E\!\left[
    B_K(Z)\bigl(Z^\top XZ-\Tr X\bigr)
  \right].
  \label{eq:app-ridge-second-Hermite-operator}
\end{align}
Use the unit-Frobenius radial direction
\begin{equation}
  X_K:=\diag\left(0,\frac{\Id_p}{\sqrt p}\right),
  \qquad
  \norm{X_K}_{\Frob}=1.
  \label{eq:app-ridge-test-direction}
\end{equation}

\begin{lemma}[Codazzi conversion and the large Hermite coupling]
\label{lem:app-ridge-Codazzi-coupling}
The vector \(\mathcal T_K[X_K]\) points in the distinguished coordinate
direction and satisfies the exact identity
\begin{equation}
  \norm{\mathcal T_K[X_K]}
  =
  \frac{s}{\sqrt p}
  \E\!\left[
    F_K''(\rho+s\Theta)\frac{R^2}{\rho}
  \right].
  \label{eq:app-ridge-Codazzi-identity}
\end{equation}
Moreover,
\begin{equation}
  \norm{\mathcal T_K[X_K]}^2
  =
  \frac K2+O(1).
  \label{eq:app-ridge-first-hermite}
\end{equation}
\end{lemma}

\begin{proof}
The constant matrix \(\alpha_K\Id\) has zero first-Hermite coefficient.
Furthermore,
\[
  \Tr(X_K\mathcal A_s)
  =
  \frac1{\sqrt p}\sum_{i=1}^p\partial_{ii}U_s
\]
is a function only of \((\Theta,R)\).  Reflection in \(Y_i\) shows that
every radial component of \(\mathcal T_K[X_K]\) vanishes.  Its distinguished
component is therefore
\[
  \frac1{\sqrt p}\sum_{i=1}^p
  \E[\Theta\,\partial_{ii}U_s(\Theta,Y)].
\]
The bounded Hessian \(\mathcal A_s=D^2U_s\) defines a tempered
distribution.  Commutation of weak derivatives, followed by Gaussian
integration by parts in the distinguished and \(i\)-th radial coordinates,
gives the Codazzi identity
\[
  \E[\Theta(\mathcal A_s)_{ii}]
  =
  \E[Y_i(\mathcal A_s)_{0i}].
\]
Indeed, writing \(\varphi\) for the standard Gaussian density,
\[
  -\int(\mathcal A_s)_{ii}\partial_0\varphi
  =
  \langle\partial_0(\mathcal A_s)_{ii},\varphi\rangle
  =
  \langle\partial_i(\mathcal A_s)_{0i},\varphi\rangle
  =
  -\int(\mathcal A_s)_{0i}\partial_i\varphi.
\]
Compactly supported cutoffs justify the pairings if desired.  From
\eqref{eq:app-ridge-ridge-Hessian},
\((\mathcal A_s)_{0i}=sF_K''(\rho+s\Theta)Y_i/\rho\).
Summation over \(i\) proves
\eqref{eq:app-ridge-Codazzi-identity}.

To estimate the expectation, put
\(q_0(\rho)=R^2/\rho=\rho-p/\rho\).  It is \(2\)-Lipschitz on
\([\sqrt p,\infty)\), and
\[
  q_0(\rho_0)=\frac{\rho_0}{2}=\frac{K^2}{4}.
\]
Using \(0\leq F_K''\leq K\) and
\(\E\abs{\rho-\rho_0}=O(1)\),
\begin{equation}
  \E[F_K''q_0(\rho)]
  =
  \overline K(s)\frac{K^2}{4}+O(K).
  \label{eq:app-ridge-coupling-expectation}
\end{equation}
Since \(p=K^4/8\), \(s^2\overline K(s)=a_K\),
\(\overline K(s)=K+O(1)\), and \(a_K=1+O(K^{-1})\),
\begin{align*}
  \norm{\mathcal T_K[X_K]}^2
  &=
  \frac{s^2}{p}
  \left(\overline K(s)\frac{K^2}{4}+O(K)\right)^2\\
  &=
  \frac{s^2\overline K(s)^2}{2}+O(K^{-1})
  =
  \frac{a_K\overline K(s)}2+O(K^{-1})
  =
  \frac K2+O(1).
\end{align*}
\end{proof}

The next identity is the polarized Hessian compatibility relation needed
for self-adjointness of the local generator.

\begin{lemma}[Second-Hermite--Bochner identity]
\label{lem:app-ridge-second-Hermite}
Let \(V\in C^2(\R^d)\), suppose \(B=\nabla^2V\) is bounded, and assume
\(\E B(Z)=\Id\).  For
\[
  \mathcal H_B[X]
  :=
  \E\!\left[B(Z)(Z^\top XZ-\Tr X)\right],
\]
one has, for every \(X,Y\in\Sym(d)\),
\begin{equation}
  \inner{Y}{\mathcal H_B[X]}_{\Frob}
  =
  \E\!\left[
    (YZ)^\top(B(Z)-\Id)(XZ)
  \right].
  \label{eq:app-ridge-second-Hermite-identity}
\end{equation}
In particular, \(\mathcal H_B\) is self-adjoint.
\end{lemma}

\begin{proof}
Let \(\varphi_d\) be the standard Gaussian density.  In the sense of
distributions,
\[
  \partial_{cd}\varphi_d=(z_cz_d-\delta_{cd})\varphi_d,
  \qquad
  \partial_{cd}B_{ab}=\partial_{bc}B_{ad},
\]
where the second equality follows from \(B=D^2V\).  Hence
\begin{equation}
  \E[(Z_cZ_d-\delta_{cd})B_{ab}(Z)]
  =
  \E[(Z_bZ_c-\delta_{bc})B_{ad}(Z)].
  \label{eq:app-ridge-Codazzi-four-index}
\end{equation}
Boundedness of \(B\) makes the distributional pairings legitimate; one may
again use compactly supported cutoffs of \(\varphi_d\).

Expanding the left side of
\eqref{eq:app-ridge-second-Hermite-identity} in indices and applying
\eqref{eq:app-ridge-Codazzi-four-index} gives
\begin{align*}
  \inner{Y}{\mathcal H_B[X]}_{\Frob}
  &=
  \sum_{a,b,c,d}Y_{ab}X_{cd}
  \E[(Z_cZ_d-\delta_{cd})B_{ab}(Z)]\\
  &=
  \E[(YZ)^\top B(Z)(XZ)]-\Tr(YX).
\end{align*}
The whitening condition gives the trace term in the second line, while
\(\Tr(YX)=\E[(YZ)^\top(XZ)]\).  This proves the identity.  Its right side is
symmetric in \(X,Y\), because \(B-\Id\) is symmetric.
\end{proof}

Set
\begin{equation}
  \tau_K^2:=\norm{\mathcal T_K[X_K]}^2,\qquad
  h_K:=\inner{X_K}{\mathcal H_K[X_K]}_{\Frob}.
  \label{eq:app-ridge-tau-h}
\end{equation}

\begin{lemma}[Second-Hermite asymptotics and the Rayleigh direction]
\label{lem:app-ridge-Rayleigh}
The test direction satisfies
\begin{equation}
  h_K=\frac K2+O(1),\qquad
  h_K-\tau_K^2=O(1).
  \label{eq:app-ridge-second-Hermite-asymptotics}
\end{equation}
Let \(L_K^{\mathrm{pack}}\) denote the operator
\begin{equation}
  L_K^{\mathrm{pack}}
  =
  \begin{pmatrix}
    \Id&\mathcal T_K/\sqrt2\\
    \mathcal T_K^\ast/\sqrt2&\Id+\mathcal H_K/2
  \end{pmatrix}
  \quad\text{on }\R^{d_K}\oplus\Sym(d_K).
  \label{eq:app-ridge-packed-generator}
\end{equation}
It is isometrically equivalent to the local generator
\eqref{eq:local-generator}.  With
\begin{equation}
  Q_K:=\frac{X_K}{\sqrt2},\qquad
  v_K:=-\frac{\mathcal T_K[X_K]}2,
  \label{eq:app-ridge-Rayleigh-vector}
\end{equation}
one has the exact quotient
\begin{equation}
  \frac{
    \inner{(v_K,Q_K)}
    {L_K^{\mathrm{pack}}(v_K,Q_K)}_{\mathrm{pack}}
  }{
    \norm{v_K}^2+\norm{Q_K}_{\Frob}^2
  }
  =
  \frac{2+h_K-\tau_K^2}{2+\tau_K^2}
  =O(K^{-1}).
  \label{eq:app-ridge-rayleigh}
\end{equation}
\end{lemma}

\begin{proof}
Since
\[
  X_KZ=\left(0,\frac{Y}{\sqrt p}\right),\qquad
  B_K-\Id=\mathcal A_s-a_K\Id,
\]
\Cref{lem:app-ridge-second-Hermite} gives
\begin{equation}
  h_K
  =
  \frac1p\E\!\left[
    F_K''\frac{R^4}{\rho^2}
  \right]
  +
  \E\!\left[
    F_K'\frac{R^2}{\rho^3}
  \right]
  -a_K.
  \label{eq:app-ridge-h-exact}
\end{equation}
For \(u=R^2/p\), the function \(u\mapsto u^2/(1+u)\) is
\(1\)-Lipschitz on \([0,\infty)\), equals \(1/2\) at \(u=1\), and
\(\norm{u-1}_{L^2}=\sqrt{2/p}=O(K^{-2})\).  Therefore
\[
  \frac1p\E\!\left[F_K''\frac{R^4}{\rho^2}\right]
  =
  \frac{\overline K(s)}2+O(K^{-1})
  =
  \frac K2+O(1).
\]
The remaining positive expectation in \eqref{eq:app-ridge-h-exact} is
bounded by
\[
  K^2\sup_{r\geq0}\frac{r^2}{(p+r^2)^{3/2}}
  \leq\frac{K^2}{\sqrt p}=2\sqrt2,
\]
and \(a_K=O(1)\).  This proves \(h_K=K/2+O(1)\).
On the other hand, the last calculation in
\Cref{lem:app-ridge-Codazzi-coupling} gives
\[
  \tau_K^2
  =
  \frac{a_K\overline K(s)}2+O(K^{-1}).
\]
Thus
\[
  h_K-\tau_K^2
  =
  \frac{(1-a_K)\overline K(s)}2+O(1)
  =
  \frac{\alpha_KK}{2}+O(1)
  =
  O(1).
\]

The map
\((v,Q)\mapsto(v,\sqrt2Q)\) is an isometry from the packed norm to
\(\norm{(v,X)}_\star^2=\norm v^2+\norm X_{\Frob}^2/2\), and conjugates
\eqref{eq:local-generator} to
\eqref{eq:app-ridge-packed-generator}.  The quadratic form of
\(L_K^{\mathrm{pack}}\) at
\eqref{eq:app-ridge-Rayleigh-vector} is
\begin{align*}
  &\norm{v_K}^2
  +2\inner{v_K}{(\mathcal T_K/\sqrt2)Q_K}
  +\norm{Q_K}_{\Frob}^2
  +\frac12\inner{Q_K}{\mathcal H_K[Q_K]}_{\Frob}\\
  &\hspace{25mm}
  =
  \frac{\tau_K^2}{4}-\frac{\tau_K^2}{2}
  +\frac12+\frac{h_K}{4}
  =
  \frac{2+h_K-\tau_K^2}{4}.
\end{align*}
The squared norm of the vector is
\((2+\tau_K^2)/4\), proving the exact quotient.  Its final order follows
from \(\tau_K^2=K/2+O(1)\) and
\eqref{eq:app-ridge-second-Hermite-asymptotics}.
\end{proof}

\begin{lemma}[Smooth realization]
\label{lem:app-ridge-mollification}
After increasing the lower threshold on \(K\), the profiles can be chosen
so that \(V_K\in C^\infty(\R^{d_K})\), while preserving exact whitening,
all asymptotic estimates above, and the common Hessian-Lipschitz constant
\(64\).
\end{lemma}

\begin{proof}
Let \(f_K(x)=(K-\abs{x-c})_+\), fix an even nonnegative
\(\eta\in C_c^\infty((-1,1))\) with \(\int\eta=1\), and choose any fixed
\(\varepsilon\in(0,1]\).  Put
\[
  \eta_\varepsilon(u)=\varepsilon^{-1}\eta(u/\varepsilon),
  \qquad
  f_{K,\varepsilon}=f_K*\eta_\varepsilon,
  \qquad
  F_{K,\varepsilon}''=f_{K,\varepsilon},
  \quad F_{K,\varepsilon}'(-\infty)=0.
\]
Convolution preserves
\[
  0\leq f_{K,\varepsilon}\leq K,\qquad
  \Lip(f_{K,\varepsilon})\leq1,\qquad
  \int_\R f_{K,\varepsilon}=K^2,
\]
and
\[
  \norm{f_{K,\varepsilon}-f_K}_{L^\infty}\leq\varepsilon.
\]
Since \(F_K''\) is \(1\)-Lipschitz and the even mollifier has zero first
moment,
\[
  F_{K,\varepsilon}'=F_K'*\eta_\varepsilon,\qquad
  \norm{F_{K,\varepsilon}'-F_K'}_{L^\infty}
  \leq\frac{\varepsilon^2}{2}.
\]
It follows uniformly for \(0\leq s\leq2K^{-1/2}\) that
\[
  \abs{\overline K_\varepsilon(s)-\overline K(s)}\leq\varepsilon,
\]
and
\[
  \abs{G_\varepsilon(s)-G(s)}
  \leq
  \frac{\varepsilon}{p}
  +\frac{\varepsilon^2}{2}
    \bigl(p^{-1/2}+p^{-3/2}\bigr).
\]
Thus \eqref{eq:app-ridge-profile-estimates} is unchanged.  The same
intermediate-value argument retunes \(s=s_{K,\varepsilon}\) and restores
exact isotropy.  Since \(f_{K,\varepsilon}(c)\geq K-\varepsilon\), the
ellipticity lower bound remains of order \(K\).  The proof of
\Cref{lem:app-ridge-Hessian-Lipschitz} uses only
\(\norm{f_{K,\varepsilon}'}_\infty\leq1\), and hence retains the same
constant \(64\).  The coupling, second-Hermite, and Rayleigh calculations
change by at most \(O(1)\), so all their asymptotic conclusions persist.
\end{proof}

\begin{proof}[Proof of \Cref{thm:local-ridge-obstruction}]
Use the smooth, retuned potential supplied by
\Cref{lem:app-ridge-mollification}.  The optimizer and whitening conditions
follow by applying the argument of
\Cref{lem:app-ridge-whitening-ellipticity} to the mollified profile, and
\begin{equation}
  \begin{gathered}
    d_K=\frac{K^4}{8}+1,\qquad
    \alphastar=\alpha_K=\Theta(K^{-1}),\\
    \betastar=\beta_K=\Theta(K),\qquad
    \kappastar=\Theta(K^2).
  \end{gathered}
  \label{eq:app-ridge-final-scaling}
\end{equation}
The common global Hessian-Lipschitz bound is
\eqref{eq:app-ridge-Hessian-Lipschitz}.

\Cref{lem:app-ridge-Codazzi-coupling} gives
\(\norm{\mathcal T_K}^2\geq K/2-O(1)\).  For the converse bound,
diagonalize an arbitrary unit-Frobenius \(X\) and write its eigenvalue
vector as \(\lambda\).  In that basis set
\[
  \widetilde T_{ji}:=\E[(B_K)_{ii}(Z)Z_j],
  \qquad
  D_{ij}:=\E[(B_K)_{ii}(Z)(Z_j^2-1)].
\]
For every \(u\), Gaussian integration by parts gives
\[
  \E[(u+\diag(\lambda)Z)^\top
      B_K(Z)(u+\diag(\lambda)Z)]
  =
  \norm u^2+2u^\top\widetilde T\lambda
  +\norm\lambda^2+\lambda^\top D\lambda.
\]
Positivity of \(B_K\) and the Schur complement yield
\(\widetilde T^\top\widetilde T\preceq\Id+D\), while
\(B_K\preceq\beta_K\Id\) gives
\(\Id+D\preceq\beta_K\Id\).  Hence
\(\norm{\mathcal T_K[X]}^2\leq\beta_K\norm X_{\Frob}^2=O(K)\).
Together with the explicit test direction,
\begin{equation}
  \norm{\mathcal T_K}^2=\Theta(K).
  \label{eq:app-ridge-operator-size}
\end{equation}

By Rayleigh--Ritz and \Cref{lem:app-ridge-Rayleigh},
\(\gstar\leq C/K\).  Conversely, the second branch in
\eqref{eq:local-Gamma} is of order at least \(K^2\), whereas
\(\betastar-1=\Theta(K)\).  Thus
\(\Gamma=\betastar-1=\Theta(K)\), and
\Cref{thm:local-spectral-sandwich} gives
\[
  \gstar
  \geq
  \max\left\{\alphastar,\frac1{4+\Gamma}\right\}
  \geq\frac cK.
\]
This proves \(\gstar=\Theta(K^{-1})\).  Substituting
\(\kappastar=\Theta(K^2)\) and
\(d_K=K^4/8+1\) yields every equivalence in
\eqref{eq:local-ridge-scaling}.
\end{proof}

\begin{remark}[Scope and minimum regularity]
\label{rem:app-ridge-scope}
The construction has
\(d_K=\Theta(\kappastar^2)\); it does not settle the fixed-dimensional
question or the sharpness of the
\(\sqrt d\log\kappastar\) branch of \(\Gamma\).
The Codazzi and second-Hermite identities use only commutation of
distributional derivatives for a bounded Hessian.  Thus the tent profile,
which is \(C^{2,1}\), already proves the analytic estimates; mollification is
used solely to realize the theorem with \(C^\infty\) potentials.
\end{remark}
